\chapter*{General Introduction}
\markboth{General Introduction}{}
\addcontentsline{toc}{chapter}{General Introduction}

Une ligne de production n'annonce pas d'ordinaire une défaillance à l'avance. Elle fonctionne, se dégrade silencieusement, puis s'arrête, et le coût n'est alors plus seulement la réparation : c'est le temps d'arrêt (\textit{downtime}), les contrôles qualité manqués, et dans les pires cas l'exposition des personnes se trouvant à proximité de la machine au moment de l'événement. Les systèmes de gestion d'actifs industriels (Enterprise Asset Management, EAM) construits autour de réparations correctives ou d'un calendrier préventif fixe n'ont jamais été conçus pour lire les données de fonctionnement propres de la machine avant ce point, ce qui constitue exactement la lacune que ce projet se propose de combler.

Cette lacune est celle autour de laquelle le stage chez SagemCom s'est construit : une plateforme EAM qui intègre la prédiction de défaillances basée sur l'apprentissage automatique (Machine Learning, ML), l'assistance des techniciens sur le terrain, et un assistant conversationnel capable de s'appuyer sur la documentation technique, le tout dans un seul et même système.

The question this report answers, concretely, is: can a single platform carry both the operational core of industrial asset management, users, machines, work orders, intervention orders, and a predictive layer capable of anticipating failures and guiding maintenance teams toward them, without the two halves working against each other?

Chapter~1 lays out the internship context, the host company, and that problem statement in full. Chapter~2 turns it into requirements. Chapter~3 designs the system that meets them. Chapter~4 walks through what was actually built. Chapter~5 tests it and reports what held up and what did not. The general conclusion returns to the question above and states plainly how much of it got answered.
